\chapter{几何与运动}

本章建立正交宇宙（\emph{Orthogonal}）中几何与运动学的基本框架。我们从洛伦兹几何（Lorentzian geometry）与黎曼几何（Riemannian geometry）的对比入手：前者以不变量 $x^{2}-t^{2}$ 刻画，后者以 $x^{2}+t^{2}$ 刻画，从而取消了类时、类空与类光方向之别。随后我们定义四维矢量的长度与点积，说明黎曼时空为何不存在任何优先的“时间方向”，进而引入能量-动量矢量与守恒定律，并以“奇异的碰撞”展示这一几何所导致的、在物理上不可回避的反直觉结论。本章进阶部分讨论不同观测者之间的能量比较、时空对称性，以及自由经典粒子的哈密顿量与拉格朗日量。

\section{洛伦兹几何与黎曼几何}

在我们居住的宇宙中，我们用来命名空间方向的词汇——上、下、左、右、前、后、南、北——全都由具体的物体来定义，诸如我们自己的身体，或我们脚下的大地。这样做只是出于方便，但归根结底我们别无选择：我们必须借助物体来辨认方向，因为空间本身的底层几何对所有方向一视同仁。

不那么明显的是，当我们谈论“未来”或“过去”时，我们实际上隐含着在指某个特定对象（或某群对象）的未来或过去。当我们说“向未来”移动一小时，我们通常指的是相对于周围的街道与建筑在空间中保持静止。但一辆疾驰的汽车在时空中划出的方向，与路边灯柱划出的方向一样，都是“未来”的一个同样合理的选择——同样地，时空的底层几何对所有这样的方向也一视同仁。

\orthofig{images/02__001.png}{洛伦兹时空中的方向}

然而，尽管你称之为“左”的方向可能正是我称之为“前”的方向，尽管当你认为我正以每秒一千公里的速度在你身旁飞掠而过时、你可能认为自己“静止不动、只在向未来前进”，但并\emph{非}时空中\emph{任意}两个方向都不可区分。归根结底，我们两人总会一致同意\emph{类时}（timelike）、\emph{类空}（spacelike）与\emph{类光}（null）方向之间的区分。

\emph{类时}方向是周围普通物体随时间划出的方向。我们或许会说一辆车正“向东行驶”，但我们当然知道它同时也在走向未来，因此在时空图上它会划出一条在走向未来时略微向东倾斜的线。

\emph{类空}方向“东”是我们可以画在图上的东西，但没有任何真实物体能沿一条\emph{只}向东、而不向时间推进的时空路径前进。

介于类时与类空方向之间的是光所划出的方向，这些被称为\emph{类光}（lightlike）或\emph{类零}（null）方向。

假设某人使用光速等于 1 的单位制（例如时间用年、距离用光年），测得时空中两个事件相隔距离 $x$、时间 $t$。那么对连接这两个事件的矢量，我们可以说：

\begin{itemize}
  \item 若 $x^{2}-t^{2}$ 为负，则该矢量是\emph{类时}的（例如图中从 A 到 B、A 到 C 的矢量）；
  \item 若 $x^{2}-t^{2}$ 为正，则该矢量是\emph{类空}的（例如图中从 A 到 E 的矢量）；
  \item 若 $x^{2}-t^{2}$ 等于零，则该矢量是\emph{类光}的（例如图中从 A 到 D 的矢量）。
\end{itemize}

狭义相对论的核心发现是：\emph{每一个}作匀速运动的观测者都会测得\emph{相同的} $x^{2}-t^{2}$ 值，因而他们会在把矢量归类为类时、类空或类光这一点上取得一致。使 $x^{2}-t^{2}$ 成为一个\emph{不变量}（invariant，即所有观测者都能认同的量）的这种时空几何，被称为\textbf{洛伦兹}几何。

\orthofig{images/02__002.png}{黎曼时空中的方向}

与之相对，在一个具有\textbf{黎曼}几何的宇宙中，所有人都认同的量是 $x^{2}+t^{2}$。这正是我们按勾股定理（Pythagoras's Theorem）在欧几里得空间中测得的两点间距离的平方。

（如果几何是弯曲的，这只有在 $x$ 与 $t$ 的测量尺度足够小时才成立，正如地球的弯曲表面只有在局限到足够小的区域时才服从欧几里得几何。）

这个平方距离 $x^{2}+t^{2}$ 恒为正（排除 $x$ 与 $t$ 同时为零的平凡情形）。于是旧的类时、类空、类光之别的消失了；在黎曼宇宙中，\emph{所有方向本质上都相同}。某个可被称为“东”的方向，与某个可被称为“未来”的方向之间，不存在任何根本差异。

在洛伦兹情形中，我们利用光速来确定 $x$ 与 $t$ 的测量单位。而在黎曼情形中，并不存在单一而普适的光速。那么我们如何选取测量距离与时间的恰当单位呢？我们只需调整时间单位与距离单位之间的比值，直到通过实验发现勾股定理始终成立——即便它同时涉及时间和空间。本注记中，我们为黎曼物理给出的所有方程，都假定已作出了这样的单位选择。

到目前为止，我们一直通过假定所关注的矢量只张成坐标系中两个方向来简化问题。而且，我们讨论过的矢量都还是\textbf{位移矢量}（displacement vector），它们像是我们可以设想的、从事件 A 指向事件 B 的箭头。本注记中我们还要讨论各种矢量，它们的分量与四个维度相关联，但度量的可能并非距离。因此一般而言，我们把矢量 $\bv$ 的分量记作 $v^{x}$、$v^{y}$、$v^{z}$ 与 $v^{t}$。（注意这些上标只是四个分量的\emph{标号}，与把某物取幂毫无关系。）四维矢量 $\bv$ 的长度平方由勾股定理从二维到四维的推广给出：

\begin{equation}
|\bv|^{2} = (v^{x})^{2} + (v^{y})^{2} + (v^{z})^{2} + (v^{t})^{2}
\end{equation}

即便观测者使用不同的坐标系、测得各个分量 $v^{x}$、$v^{y}$、$v^{z}$ 与 $v^{t}$ 的不同数值，所有人仍会认同矢量 $\bv$ 的长度 $|\bv|$。

\section{点积}

在黎曼几何中，使用不同坐标系的观测者除了认同矢量的长度，还会认同矢量之间的\emph{夹角}。这使我们能够定义一个同时与长度和夹角都相关的、不变量般的量。

设有两个矢量 $\bv$ 与 $\bw$，我们如下定义 $\bv$ 与 $\bw$ 的\textbf{点积}（dot product）：对四个维度 $x$、$y$、$z$、$t$ 中的每一个，把两个矢量的分量相乘，再全部相加。我们把它记作 $\bv\cdot\bw$，于是：

\begin{equation}
\bv\cdot\bw = v^{x}w^{x} + v^{y}w^{y} + v^{z}w^{z} + v^{t}w^{t}
\tag{黎曼点积}
\end{equation}

按此定义，矢量的长度平方就是它与自身的点积：

\begin{equation}
|\bv|^{2} = \bv\cdot\bv
\end{equation}

点积在旋转下保持不变。为看出这一点，考虑一个三角形，其中两条边是矢量 $\bv$ 与 $\bw$；第三条边将是 $\bv-\bw$。第三条边的长度平方为：

\begin{equation}
|\bv-\bw|^{2} = (\bv-\bw)\cdot(\bv-\bw) = |\bv|^{2} + |\bw|^{2} - 2\,\bv\cdot\bw
\end{equation}

因此 $\bv\cdot\bw = \bigl(|\bv|^{2} + |\bw|^{2} - |\bv-\bw|^{2}\bigr)/2$。由于旋转三角形时这些长度都不变，点积也不可能改变。

\orthofig{images/02__003.png}{点积}

若 $\bv$ 与 $\bw$ 之间的夹角为 $\theta$，则：

\begin{equation}
\bv\cdot\bw = |\bv|\,|\bw|\cos\theta
\end{equation}

这并非一个显而易见的论断，但很容易证明。设想最简单的情形：$\bw$ 位于 $x$ 轴正半轴上（于是 $w^{y}=w^{z}=w^{t}=0$，且 $w^{x}>0$），而 $\bv$ 位于 $xt$ 平面内（于是 $v^{y}=v^{z}=0$）。由于众多分量为零，我们有：

\begin{equation}
\bv\cdot\bw = v^{x}w^{x}, \qquad |\bw| = w^{x}
\end{equation}

由余弦的定义，我们又有：

\begin{equation}
v^{x} = |\bv|\cos\theta
\end{equation}

由此可知，在这一特定情形下：

\begin{equation}
\bv\cdot\bw = |\bv|\,|\bw|\cos\theta
\end{equation}

但对于任意给定的 $\bv$ 与 $\bw$，我们总可以旋转它们，使其满足以上所设的限制条件。旋转不会改变矢量的长度、也不改变它们之间的夹角 $\theta$，更不会改变它们的点积；因此若在旋转后结论成立，那么它在旋转之前也必然一直成立。

夹角的余弦在夹角小于 $90^\circ$ 时为正，恰为 $90^\circ$ 时为零，介于 $90^\circ$ 与 $180^\circ$ 之间时为负。于是：

\begin{itemize}
  \item 若 $\bv\cdot\bw > 0$，则 $\bv$ 与 $\bw$ 大致指向\emph{同一}方向（夹角小于 $90^\circ$）；
  \item 若 $\bv\cdot\bw = 0$，则 $\bv$ 与 $\bw$ 相互垂直；
  \item 若 $\bv\cdot\bw < 0$，则 $\bv$ 与 $\bw$ 大致指向\emph{相反}方向（夹角介于 $90^\circ$ 与 $180^\circ$ 之间）。
\end{itemize}

任一矢量 $\bv$ 都可以沿某个单位矢量 $\bu$ 分解为平行与垂直两个部分：

\begin{equation}
\bv = (\bu\cdot\bv)\,\bu + \bigl(\bv - (\bu\cdot\bv)\,\bu\bigr)
\end{equation}

第一部分显然平行于 $\bu$，因为它不过是数值 $\bu\cdot\bv$ 乘以 $\bu$。点积 $\bu\cdot\bv$ \emph{几乎}就是 $\bv$ 在 $\bu$ 上的投影长度；唯一的差别在于，若两矢量大致反向，则 $\bu\cdot\bv$ 为负，此时它是投影长度的相反数。

第二部分与 $\bu$ 垂直，这可以通过它与 $\bu$ 取点积看出：

\begin{equation}
\bu\cdot\bigl(\bv - (\bu\cdot\bv)\,\bu\bigr)
= (\bu\cdot\bv) - (\bu\cdot\bv)(\bu\cdot\bu)
= (\bu\cdot\bv) - (\bu\cdot\bv)
= 0
\end{equation}

我们将沿坐标系四个轴的单位矢量称为 $\be_{x}$、$\be_{y}$、$\be_{z}$、$\be_{t}$；这四个矢量被称为该坐标系的\textbf{基矢量}（basis vector）。于是对任意矢量 $\bv$，它在该坐标系中的分量由它与基矢量的点积给出：

\begin{equation}
v^{i} = \be_{i}\cdot\bv, \qquad i = x,\,y,\,z,\,t
\end{equation}

在黎曼宇宙中，所有这些矢量几何都非常简单：它不过是我们所熟悉的欧几里得空间几何再加上时间这一维。与之相反，在洛伦兹宇宙中，我们必须始终小心时间维度上的符号差异。例如，在洛伦兹宇宙中，两个四维矢量（或“四矢量”）的点积对时间坐标带有负号：

\begin{equation}
\bv\cdot\bw = v^{x}w^{x} + v^{y}w^{y} + v^{z}w^{z} - v^{t}w^{t}
\tag{洛伦兹点积}
\end{equation}

\section{没有类时方向的时间}

一个其几何中并未内蕴任何时间概念的宇宙，可能会让人联想到一幅类似我们宇宙中某次快照的景象：某一个单一的、冻结而静止的瞬间。但稍加思索便会发现，这完全说不通。当我们从全部四个维度来看时，我们自己的宇宙同样也是不变的；对我们而言构成\emph{变化}的，是穿过四维宇宙的、一个个相继的三维切片彼此不同这一事实。因此，如果支配黎曼宇宙内容的物理定律也允许相继的三维切片彼此不同，那么产生一种被感知为\emph{变化}的东西的可能性就同样存在。

\orthofig{images/02__005.png}{由世界线定义时间}

相对论告诉我们，并不存在唯一“正确的”方式把我们自己的宇宙切分为一个个相继的时间瞬间。处于相对运动中的不同人会以不同的方式做这种切分，每个人都把自己划出的\textbf{世界线}（world line，即他在时空中留下的路径或历史）当作时间方向的标记，而把与其垂直的所有方向当作空间方向。

现在，一个黎曼宇宙究竟包含何种物体或许并不显然，但结果它也会充满世界线。（其深层原因在于量子力学，但当我们接下来探讨波的时候，会一窥世界线是如何浮现的。）主要差别在于：洛伦兹宇宙中的世界线\emph{必须}沿类时方向延伸，而在黎曼宇宙中，它们可以沿的方向不受任何限制。

我们通常把时间看作普适的，而非我们自己世界线的一种属性，这是因为日常生活中我们周围一切人与物彼此之间的速度都非常小，于是我们所有的世界线都挤在一个紧密的线束中一起延伸，只是来回微小地穿插，但基本上互相平行。正是\emph{这一点}，在大多数实际情形下赋予了我们对时间的单一而共同的概念。对日常经验而言最重要的，是这种共享的运动，而非底层时空几何把某一类（大得多的！）方向归类为类时这一事实。黎曼宇宙中一个类似的、近乎平行的世界线线束，也会产生同样的实际感受，使线束的方向成为一种有用的共享时间概念。

因此，我们不妨假定自己能想象自身是黎曼宇宙的居民，拥有自己的世界线。正如我们在自己宇宙中所做的那样，我们通常会采用一个时间坐标，其轴指向我们世界线的方向，而空间坐标的轴与其正交。我们还将假定，我们有一套明确的方式，沿着我们的世界线选定某个方向作为“未来”。与我们自己的宇宙一样，这归根结底是熵的问题：未来将是熵增大的那个方向。

任何世界线与我们的世界线平行的物体，我们都会认为它是静止的，因为在我们坐标系中它的空间位置不随时间改变。任何世界线与我们的世界线成一定倾角的物体，我们都会认为它在运动。

\orthofig{images/02__006.png}{四维速度与四维加速度}

我们有自己的时间坐标 $t$，但相对论要求任何物体的时间流逝都沿着它的世界线来度量。我们把这样测得的时间称为\textbf{固有时}（proper time）$\tau$。在黎曼宇宙中，两个事件之间对某物体流逝的固有时，就是该物体世界线的长度——除了单位选择的问题之外，它与沿任何其他曲线测得的长度相同。

现在，考虑一个与物体世界线相切的单位矢量 $\bu$。我们称它为物体的\textbf{四维速度}（four-velocity），因为它是一个描述物体瞬时运动状态的四维矢量。我们可以通过用 $\bu$ 的空间分量除以其时间分量，来提取物体在某个特定坐标系中的普通（三维）速度 $\bv$；在图中，一切都被限制在 $xt$ 平面内，于是我们有：

\begin{equation}
v^{x} = u^{x}/u^{t}
\end{equation}

因为如果物体沿 $\bu$ 所标示的切线继续运动，它将在时间 $u^{t}$ 内行进距离 $u^{x}$。

当然，图中的物体并未保持恒定速度，它在加速。物体的\textbf{四维加速度}（four-acceleration），我们记作 $\mathbf{a}$，定义为物体的四维速度关于固有时 $\tau$ 的变化率。我们写作：

\begin{equation}
\mathbf{a} = \partial_{\tau}\bu
\end{equation}

不必被这里的记号吓到！符号“$\partial$”不过是“变化率”的一种简洁写法，而下标 $\tau$ 意味着我们谈论的是当 $\tau$ 变化时 $\bu$ 所发生的变化。

由于四维速度 $\bu$ 始终是单位矢量，它唯一可能发生的变化是方向上的，而非长度上的，这意味着四维加速度 $\mathbf{a}$ 必须与 $\bu$ 正交。

\orthofig{images/02__007.png}{恒定加速度}

假设我们取一个在我们的坐标系中初始静止的物体，并让它沿我们的 $x$ 方向受到恒定加速。确切地说，我们要求物体加速度的\emph{大小}为常量 $a$，且加速度四矢量始终留在 $xt$ 平面内。物体的世界线将起初平行于我们的 $t$ 轴，然后逐渐向 $x$ 轴倾斜。

在洛伦兹宇宙中，物体的运动会越来越快，但它永远达不到光速。

在牛顿宇宙中，物体的速度没有上限，但在任何有限时刻其速度仍是有限的。

在黎曼宇宙中，物体的世界线将是一段圆弧。这并不应令人惊讶，因为这里涉及的几何与数学非常熟悉——它与离心力的数学大致相同：对一个运动物体施以恒定的侧向拉力，会使它沿圆周路径运动。

我们可以把物体的 $(x,t)$ 坐标、四维速度与四维加速度显式地写成固有时 $\tau$ 的函数：

\begin{align}
x(\tau) &= \tfrac{1}{a}\,(1-\cos a\tau,\ \sin a\tau),\\
\bu(\tau) &= \partial_{\tau}x(\tau) = (\sin a\tau,\ \cos a\tau),\\
\mathbf{a}(\tau) &= \partial_{\tau}\bu(\tau) = a\,(\cos a\tau,\ -\sin a\tau).
\end{align}

于是，在有限时间之后——按我们的坐标是 $1/a$，按物体自身是 $\pi/(2a)$——物体的四维速度将平行于我们的 $x$ 轴，我们便会认为它具有无限大的速度。在那之后，它将在我们的坐标中反向于时间运动！

至于这种恒定加速在黎曼宇宙中\emph{实际上}究竟能在何种条件下维持，将取决于本应施加于物体的力的机制；与我们宇宙中一样，任何这样的过程都会受制于关于能量与熵的严格定律。但作为纯运动学的陈述，图示给出了\emph{如果}恒定加速能够实现时会发生什么。

\section{能量-动量矢量}

假设我们给周围每一个物体都关联一个箭头，其长度等于物体的质量 $m$，方向与物体的世界线相切。我们总选取这些箭头指向与我们个人的“未来”选择大致相符的方向；确切地说，我们总选取它们的指向，使其与我们的时间轴所成的角不超过 $90^\circ$。

\orthofig{images/02__008.png}{能量-动量矢量}

这个箭头就是物体的\textbf{能量-动量矢量}（energy-momentum vector），我们用符号 $\bP$ 表示它。我们已经定义了与物体世界线相切的\emph{单位}矢量，即四维速度 $\bu$，因此可以写出：

\begin{equation}
\bP = m\,\bu
\end{equation}

为简单起见，我们暂且假设物体世界线位于我们的 $xt$ 平面内，于是只需处理两个维度。我们把能量-动量矢量的 $t$ 分量称为物体的\textbf{总能量}（total energy）$E$，把它的 $x$ 分量称为物体的\textbf{动量}（momentum）$p$。于是：

\begin{equation}
P^{t} = E, \qquad P^{x} = p
\end{equation}

\begin{equation}
m^{2} = |\bP|^{2} = (P^{t})^{2} + (P^{x})^{2} = E^{2} + p^{2}
\end{equation}

物体的速度 $v$ 将是其四维速度的 $x$ 分量与 $t$ 分量之比 $u^{x}/u^{t}$。由于能量-动量矢量是四维速度的倍数，我们可以直接从能量-动量矢量的分量得到这个比值：

\begin{equation}
v = p/E
\end{equation}

利用上面最后两个方程，我们有：

\begin{align}
p &= Ev,\\
E^{2} &= m^{2} - p^{2} = m^{2} - E^{2}v^{2},\\
E^{2}(1 + v^{2}) &= m^{2},
\end{align}

从而给出用 $m$ 与 $v$ 表示的 $E$ 与 $p$：

\begin{equation}
E = \frac{m}{\sqrt{1+v^{2}}}, \qquad p = \frac{mv}{\sqrt{1+v^{2}}}
\end{equation}

你可能会认出，这些公式是对通常的相对论公式的一个简单改动，其中 $1+v^{2}$ 取代了 $1-v^{2}$。

利用 $1/\sqrt{1+v^{2}}$ 的泰勒级数，我们可以把这两个公式在 $v$ 较小时近似为：

\begin{equation}
\frac{1}{\sqrt{1+v^{2}}} \approx 1 - \tfrac{1}{2}v^{2},
\qquad
E \approx m - \tfrac{1}{2}mv^{2},
\qquad
p \approx mv
\end{equation}

\orthofig{images/02__009.png}{牛顿、洛伦兹与黎曼物理中的动能}

这里动量的近似值 $mv$，正是普通牛顿物理中物体的动量。而在牛顿物理中，物体的动能为：

\begin{equation}
K_{\text{Newtonian}} = \tfrac{1}{2}mv^{2}
\end{equation}

因此总能量的近似值为：

\begin{equation}
E \approx m - K_{\text{Newtonian}}
\end{equation}

在我们自己的宇宙中，我们认为物体具有一份静质量能量 $mc^{2}$（在 $c=1$ 的单位下等于 $m$），把物体的动能\emph{加}到它上面，即得总能量。而在黎曼宇宙中，我们必须从静质量能量中\emph{减去}动能。我们把精确的黎曼动能 $K$ 定义为：

\begin{equation}
K = m - E = m\left(1 - \frac{1}{\sqrt{1+v^{2}}}\right)
\end{equation}

在牛顿物理与洛伦兹物理中，随着物体速度增大，其动能可以无界增长。在牛顿物理中速度没有上限，因而 $K_{\text{Newtonian}}$ 可以任意大。在洛伦兹物理中

\begin{equation}
K_{\text{Lorentzian}} = m\left(\frac{1}{\sqrt{1-v^{2}}} - 1\right)
\end{equation}

也再次可以任意大。而在黎曼物理中，与牛顿物理一样，你能测得的物体速度没有上限；若物体的世界线垂直于你的世界线，其速度将是无穷大。但无论速度多大，其动能永远不会超过它的静质量 $m$。

那么势能呢？我们可以用通常的方式、借助作用在感兴趣粒子上的任何力来定义势能，于是动能与势能之和将被守恒。稍后我们将深入探讨黎曼版电磁学，看看一种静电势能将呈现何种形式。但我们可以立刻注意到一个奇特的事实。两个粒子彼此束缚所组成的系统，其势能通常比它们自由时更低，因为在束缚态中粒子“被困在一个能量势阱里”，四周是向下倾斜的阱壁。在洛伦兹物理中，这导致“质量亏损”：束缚系统的质量小于其各组成部分之和。在黎曼物理中，由于动能与势能相对于静质量能量具有相反的符号意义，质量亏损将朝相反的方向起作用：束缚系统的质量将\emph{大于}其各组成部分之和。

对于小速度，黎曼物理看上去与牛顿物理别无二致，动能“倒过来”这一事实未必显现。当固定数目的粒子只是彼此弹开时，动能不过是从一个粒子传递给另一个，或转化为势能。但在任何可以产生新粒子的情形中，动能与静质量能量的这种关系会产生惊人的后果：新粒子的产生不是由产生它的系统的动能或势能的\emph{减少}来平衡，而是由\emph{增加}来平衡。我们已经指出，在黎曼物理中，束缚系统的质量大于其各组成部分，因此任何处于束缚态、且能够（比如说）发光的物质，都将面临把它多余的质量转化为光、并在此过程中把自身瓦解的风险！这些后果——以及物质何以能在这种条件下保持稳定这一谜题——在小说中有详细得多的探讨。

\section{守恒定律}

在定义了能量-动量矢量之后，我们应当能够说明，在黎曼宇宙中，能量与动量被某个物体间的相互作用\emph{守恒}意味着什么。我们可以坚持选定一个特定的参考系，在其中给所有相关物体赋予能量-动量矢量——并确保它们都指向该参考系的未来——然后说，如果各物体在相互作用\emph{之前}所有能量-动量矢量之和等于相互作用\emph{之后}之和，那么就满足守恒定律。

\orthofig{images/02__010.png}{能量-动量守恒}

然而，还有一种优雅得多、完全无须援引任何特定参考系的做法。我们只需沿每个物体的世界线画一个箭头，令其长度等于该物体的质量，并选取所有箭头都\emph{背离}相互作用点。如果能量与动量守恒，那么所有这些箭头之和将为零。

确立了守恒律的这种对称表述之后，我们就能看出，它对任何一位真的前来、为相互作用赋予“之前”与“之后”能量-动量矢量的观测者意味着什么。

\orthofig{images/02__011.png}{能量-动量守恒}

原先的箭头全部背离相互作用点，因此观测者认为是与相互作用\emph{之前}的物体相关联的那些箭头，都需要被反向，以使它们指向该观测者的未来。但由于原先所有箭头之和为零，这些被反向的箭头之和将等于其余所有箭头之和。于是对这位观测者而言，无论他们把哪个方向称为“未来”，能量与动量在相互作用前后都是相同的。

观测者该如何处理一条\emph{恰好}垂直于其选定时间轴的世界线呢？在这种情况下，事实上唯一能做的就是任意地决定：要么把所有这类物体都视为从该相互作用中产生，此时它们的箭头保持不变，被当作相互作用\emph{之后}的能量-动量；要么把所有这类物体都视为抵达该相互作用，此时它们的箭头被反向，被当作相互作用\emph{之前}就已存在的能量-动量。无论哪种方式，能量与动量都将被看到是守恒的。

\section{奇异的碰撞}

\orthofig{images/02__012.png}{碰撞一}

\orthofig{images/02__013.png}{碰撞二}

假设黎曼宇宙中两个点状物体发生一次正碰。我们假定这是一次\textbf{完全弹性}（perfectly elastic）碰撞，没有能量转化为其他形式，两个物体都不碎裂，而且两者不会粘在一起。

在这些假定下，如果我们知道参与碰撞的每个粒子的一个能量-动量矢量，那么利用守恒定律就能求出其余矢量。例如在\textbf{右上方}的图中，如果我们知道第一个粒子的\emph{初始}能量-动量矢量 $\bP_{i,1}$ 与第二个粒子的 $\bP_{i,2}$，我们就能算出它们的\emph{末态}能量-动量矢量 $\bP_{f,1}$ 与 $\bP_{f,2}$ 必须是什么。在此图所示的碰撞中，能量与动量守恒，于是在黎曼宇宙中任何安排两个粒子以这种方式碰撞的人，都会预期结果大致如这里所描绘的那样。

在\textbf{右下方}的图中，我们重画了与之前完全相同的碰撞，只不过它被旋转了大约 $60^\circ$，并且我们不得不交换 $\bP_{f,2}$ 箭头的方向，以使它在新的参考系中具有正能量。这仍然是一次完全有效的碰撞，服从守恒定律……但对该图所示时间轴上的观测者而言，现在是三条世界线汇聚到一起，碰撞之后只留下一条。在这种观点下，第二个物体与自身的反物质版本相撞，两者湮灭，释放出的能量被第一个物体吸收，使其速度大大降低（总能量与动能的符号意义相反）。

两位观测者对同一组事件给出不同描述，这在相对论中是常有的事，但真正的转折出现在我们设想某个人\emph{安排}这样一次碰撞意味着什么的时候。在第一种观点中，我们可以说，任何按图示初始能量-动量矢量搭建实验的人，都能用守恒定律预测结果。但在第二种观点中，同一张图却告诉我们：任何成功安排两个粒子以能量-动量矢量 $\bP_{i,1}$ 与 $\bP_{i,2}$ 相撞的人，只有当这次成功伴随着第二个粒子的反物质版本在碰撞\emph{之前}抵达、一同参与碰撞时，才会成功。除非第三种成分也在场，否则没有人能让最初两个粒子以这些特定的速度相撞。

基于在我们自己宇宙中的经验，我们的直觉可能会反驳说，理应存在一种实验，其中只有三个粒子中的两个参与碰撞。但守恒定律说这是\emph{物理上不可能的}。在我们自己的宇宙中，我们习惯于物理定律约束我们能做什么，但从不是以这种方式！

\begin{keybox}[本节的结论]
这是本书最有冲击力的结论之一：在黎曼宇宙中，能量-动量守恒不仅告诉我们“会发生什么”，还会反过来限制“什么能够被安排出来”。同一张时空图存在两种自洽解读——两个粒子弹性散射，或者两个粒子中一个与自身的反物质湮灭——二者都满足守恒定律。因此，某个观测者认为只需两颗粒子的碰撞过程，在另一个观测者看来必然同时伴随着第三个粒子（反物质）的出现；两颗粒子以该方式、该速度碰撞这件事，在物理上不可能单独发生。
\end{keybox}

关于这一主题，尚有给不介意稍高数学水平的读者阅读的进阶材料。

\section{进阶：不同观测者之间的能量比较}

在几何与运动的主页中，我们把物体的能量-动量矢量定义为长度等于物体质量、指向沿物体世界线的矢量。当然，沿任一条线都有\emph{两个}方向，因此我们给出的规则是：选取与描述该物体的观测者的时间轴所成角不超过 $90^\circ$ 的那个方向。

当你致力于在单一参考系中描述一切时，这是一套有用的方案；但如果你有两位处于相对运动的观测者，就会出现一些情形，使得这条方向规则令他们沿物体的世界线选取\emph{相反}的方向，于是两种描述将不相容。

在洛伦兹宇宙中，这个问题不会出现，因为处于相对运动的普通观测者仍然认同任何矢量时间分量的\emph{符号}。我们当然可以想象时间反演的观测者，但他们构成一个独特的类别，对我们而言其\emph{所有}时间分量都取负。但在黎曼宇宙中，如果你在一组速度各异的物体的一批世界线上放置指向你正向时间的箭头，而我以相对于你相当高的速度运动，那么我可能在有些情形下认同你的箭头，在另一些情形下则不认同。如果我们各自把所有能量-动量矢量相加，所得的两个总矢量将毫无关联。我们当然在用不同的坐标系，但那并非问题所在；把一个矢量从一个坐标系变换到另一个坐标系是件直截了当的事。但若把这些撇开，仅把各个能量-动量矢量 $\bP_{i}$ 看作坐标无关的几何对象，你的总和可能是 $\bP_{1}+\bP_{2}+\bP_{3}+\bP_{4}$，而我的却是 $\bP_{1}+\bP_{2}-\bP_{3}-\bP_{4}$，于是它们是完全不同的矢量。

由于沿世界线选取哪个方向确实没有\emph{物理依据}\footnote{如果物体带有电荷，沿其世界线的两个方向就不再不可区分。这一点在黎曼电磁学一节中进一步讨论。}，四维空间中的能量-动量矢量在符号上是不确定的。好消息是，物体的\textbf{应力-能量张量}（stress-energy tensor）——它封装了关于能量与动量的密度和流动的信息——并不依赖于沿世界线选择方向。一个在其静止系中密度为 $\rho$ 的物体的应力-能量张量的坐标矩阵为：

\begin{equation}
T^{ab} = \rho\, u^{a} u^{b}
\end{equation}

其中 $\bu$ 是物体的四维速度。显然，在这里交换 $\bu$ 的符号毫无影响！因此，若一位观测者在其坐标系中计算应力-能量张量，另一位观测者在其自己的坐标系中计算它，两者将通过张量变换的正常几何规则相互关联。

于是，能量-动量守恒体现为：应力-能量张量（对一切事物，包括所有物质与辐射）的散度为零。

话虽如此，选定一个参考系并在该参考系中计算系统的总能量与总动量，仍然没有任何\emph{错}。我们可以在相对论性的哈密顿形式体系中做物理——下文将加以描述——其中一旦选定了时间坐标，我们就着手定义一个作为哈密顿量的正能量。在该形式体系中，总能量与总动量将是守恒量。唯一的错误在于，把这些守恒量合并成一个四维矢量，并指望该四维矢量能通过一次旋转与别人以同样方式得到的矢量相关联。

\section{进阶：对称性}

为简单起见，假设一个四维黎曼宇宙在一切方向上无限且完全平坦。这是一种理想化，但它非常有用；我们自己的宇宙有一种类似的理想化，称为闵可夫斯基时空（Minkowski spacetime），它是狭义相对论和大多数相对论量子物理的基础。我们把这个理想化的宇宙称为\textbf{欧几里得}四维空间，因为它服从欧几里得几何的公设。

[\textbf{注意}：在物理学文献中，“欧几里得”一词常被用于指经由 Wick 转动（Wick rotation）所产生的那些版本的物理定律。那些\emph{不是}适用于\emph{Orthogonal}宇宙的定律。]

我们选取一个事件 O 作为坐标原点，选取四个互相垂直的方向作为 $x$、$y$、$z$、$t$ 轴，并选定一套单位制。于是该宇宙中任何事件都可用一组四元实数 $(x,y,z,t)$ 来指定。所有这样的四元组的集合，数学家称之为 $\mathbb{R}^{4}$。

$\mathbb{R}^{4}$ 是一个矢量空间，对于 $\mathbb{R}^{4}$ 中任何矢量 $\bv=(v^{x},v^{y},v^{z},v^{t})$，我们用标准点积求其长度 $|\bv|$：

\begin{equation}
|\bv|^{2} = \bv\cdot\bv = (v^{x})^{2} + (v^{y})^{2} + (v^{z})^{2} + (v^{t})^{2}
\end{equation}

在一般的弯曲宇宙中，我们不能对两个事件的坐标做加减，但在这个平坦宇宙中，这完全说得通。从事件 $\mathbf{a}=(a^{x},a^{y},a^{z},a^{t})$ 到事件 $\mathbf{b}=(b^{x},b^{y},b^{z},b^{t})$ 的矢量就是 $\mathbf{b}-\mathbf{a}=(b^{x}-a^{x},\,b^{y}-a^{y},\,b^{z}-a^{z},\,b^{t}-a^{t})$，而两事件之间的距离为 $|\mathbf{b}-\mathbf{a}|$。

现在假设我们的欧几里得四维空间中充满了各种物体，它们满足本质上属于几何性质的定律（一个例子就是能量-动量守恒定律）。通过考察该空间的\textbf{对称性}（symmetry）集合——即那些改变物体位置或取向、却使其内禀几何性质（诸如两事件间的距离、两条世界线间的夹角、或两个波前之间的间距）保持不变的运算——我们可以更深入地理解这类几何定律可能是什么。

形式上，我们把函数 $f:\mathbb{R}^{4}\to\mathbb{R}^{4}$ 称为一个对称性，是指对任意两个事件 $\mathbf{a}$ 与 $\mathbf{b}$，$f(\mathbf{a})$ 与 $f(\mathbf{b})$ 之间的距离等于 $\mathbf{a}$ 与 $\mathbf{b}$ 之间的距离：

\begin{equation}
|f(\mathbf{b})-f(\mathbf{a})| = |\mathbf{b}-\mathbf{a}|
\end{equation}

直观上不难想到三类满足此条件的函数，用几何语言描述为：

\begin{itemize}
  \item \emph{平移}（translation）：$f$ 把一切沿固定方向移动固定距离；
  \item \emph{旋转}（rotation）：$f$ 至少固定一个事件，并刚性地把它周围的一切绕之转动；
  \item \emph{反射}（reflection）：反转图形的取向，同时保持距离与夹角不变。
\end{itemize}

平移很容易用数学描述。若 $\mathbf{s}$ 是任意矢量，我们可以定义：

\begin{equation}
T_{\mathbf{s}}(\mathbf{a}) = \mathbf{a} + \mathbf{s}
\end{equation}

$T_{\mathbf{s}}$ 显然是对称性，因为：

\begin{equation}
T_{\mathbf{s}}(\mathbf{b}) - T_{\mathbf{s}}(\mathbf{a}) = \mathbf{b}-\mathbf{a}
\end{equation}

由于平移可以把坐标原点移到任何其他事件，我们将专注于那些\emph{固定原点}的旋转与反射。它们总能与平移组合，给出最一般的对称性。

一个固定原点的对称性 $R$ 将是 $\mathbb{R}^{4}$ 上的线性函数，这意味着我们可以在施加它之前或之后对矢量做加减、数乘，结果相同：

\begin{equation}
R(\bv+\bw) = R(\bv) + R(\bw), \qquad R(s\bv) = s\,R(\bv)
\end{equation}

与 $\mathbb{R}^{4}$ 上任何线性函数相伴的是一个 $4\times4$ 矩阵 $R^{i}{}_{j}$，满足：

\begin{equation}
R(\bv) = R(v^{j}\be_{j}) = R^{i}{}_{j}\,v^{j}\,\be_{i}
\end{equation}

其中矢量分量 $v^{j}$ 与矩阵分量 $R^{i}{}_{j}$ 都是就标准基 $\{\be_{x},\be_{y},\be_{z},\be_{t}\}$ 而言的，并且我们使用了爱因斯坦求和约定。换言之，我们可以用矩阵去乘矢量的分量，从而再现 $\mathbb{R}^{4}$ 这一矢量空间上的任何线性函数。

若 $R$ 是一个固定原点的对称性，它将保持 $\mathbb{R}^{4}$ 上的标准点积不变，因此对任意两个矢量 $\bv$ 与 $\bw$：

\begin{equation}
R(\bv)\cdot R(\bw) = \bv\cdot\bw
\end{equation}

这不过是说 $R$ 将保持矢量的长度以及它们之间的夹角不变，因为这些量都依赖于点积。

令 $\bv=\be_{i}$、$\bw=\be_{j}$，其中 $\be_{i}$ 是 $\mathbb{R}^{4}$ 上标准基的第 $i$ 个基矢量，即一个第 $i$ 分量为 1、其余分量皆为零的四元组。于是我们有：

\begin{equation}
R(\be_{i})\cdot R(\be_{j}) = \be_{i}\cdot\be_{j}
\end{equation}

标准基是正交归一的，故 $\be_{i}\cdot\be_{j}=\delta_{ij}$（克罗内克 $\delta$ 符号，$i=j$ 时为 1、$i\neq j$ 时为 0）……这也正是 $4\times4$ 单位矩阵 $I_{4}$ 的 $i,j$ 分量：

\begin{equation}
R(\be_{i})\cdot R(\be_{j}) = \be_{i}\cdot\be_{j} = (I_{4})^{i}{}_{j}
\end{equation}

更有：

\begin{equation}
R(\be_{i})\cdot R(\be_{j}) = R^{k}{}_{i}R^{k}{}_{j} = (R^{\!\top})^{i}{}_{k}R^{k}{}_{j} = (R^{\!\top}R)^{i}{}_{j}
\end{equation}

其中 $R^{\!\top}$ 是 $R$ 的转置，即把 $R$ 的行变为列所得的矩阵。由于 $R^{\!\top}R$ 与 $I_{4}$ 的所有分量都一致，两者必相等：

\begin{equation}
R^{\!\top}R = I_{4}
\end{equation}

于是 $R^{\!\top}$ 就是 $R$ 的逆。

这给了我们一种刻画欧几里得四维空间所有线性对称性的方法：这类函数矩阵的\emph{转置}同时也是它的\emph{逆}。所有这类 $4\times4$ 矩阵的集合称为 $\mathrm{O}(4)$，即四维\textbf{正交群}（orthogonal group）。

举几个具体的例子，下面这个矩阵：

\begin{equation}
\begin{pmatrix}
\cos\theta & -\sin\theta & 0 & 0\\
\sin\theta & \cos\theta & 0 & 0\\
0 & 0 & \cos\phi & -\sin\phi\\
0 & 0 & \sin\phi & \cos\phi
\end{pmatrix}
\end{equation}

在 $xy$ 平面内旋转角度 $\theta$，在 $zt$ 平面内旋转角度 $\phi$。把它乘以它的转置，得到：

\begin{equation}
\begin{pmatrix}
\cos\theta & \sin\theta & 0 & 0\\
-\sin\theta & \cos\theta & 0 & 0\\
0 & 0 & \cos\phi & \sin\phi\\
0 & 0 & -\sin\phi & \cos\phi
\end{pmatrix}
\begin{pmatrix}
\cos\theta & -\sin\theta & 0 & 0\\
\sin\theta & \cos\theta & 0 & 0\\
0 & 0 & \cos\phi & -\sin\phi\\
0 & 0 & \sin\phi & \cos\phi
\end{pmatrix}
=
\begin{pmatrix}
1 & 0 & 0 & 0\\
0 & 1 & 0 & 0\\
0 & 0 & 1 & 0\\
0 & 0 & 0 & 1
\end{pmatrix}
\end{equation}

一个更简单的例子是下面这个矩阵：

\begin{equation}
\begin{pmatrix}
-1 & 0 & 0 & 0\\
0 & 1 & 0 & 0\\
0 & 0 & 1 & 0\\
0 & 0 & 0 & 1
\end{pmatrix}
\end{equation}

它是一个反转 $x$ 坐标的反射。把它乘以它的转置（也就是同一个矩阵），得到：

\begin{equation}
\begin{pmatrix}
-1 & 0 & 0 & 0\\
0 & 1 & 0 & 0\\
0 & 0 & 1 & 0\\
0 & 0 & 0 & 1
\end{pmatrix}
\begin{pmatrix}
-1 & 0 & 0 & 0\\
0 & 1 & 0 & 0\\
0 & 0 & 1 & 0\\
0 & 0 & 0 & 1
\end{pmatrix}
=
\begin{pmatrix}
1 & 0 & 0 & 0\\
0 & 1 & 0 & 0\\
0 & 0 & 1 & 0\\
0 & 0 & 0 & 1
\end{pmatrix}
\end{equation}

我们可以通过取矩阵的行列式来区分纯旋转与反射。旋转的行列式为 1，反射的行列式为 $-1$。所有转置为其逆\emph{且}行列式为 1 的 $4\times4$ 矩阵的集合称为 $\mathrm{SO}(4)$，即四维\textbf{特殊正交群}（special orthogonal group）。

把正交群 $\mathrm{O}(4)$ 的任意元素与平移组合所得的变换群，称为四维\textbf{欧几里得群}（Euclidean group）$\mathrm{E}(4)$。这个群给出四维欧几里得空间的所有对称性。若我们只限于旋转与平移——略去反射——我们就得到 $\mathrm{SE}(4)$，即\textbf{特殊欧几里得群}（special Euclidean group）。

把包含平移的欧几里得群的一般对称性用矩阵来表示有时会很方便。我们可以采用这样的约定：在 $\mathbb{R}^{4}$ 中一个事件的四个坐标末尾再补上一个坐标 1。按这种记法，我们可以把一个线性运算 $R$ 与一个平移矢量 $\mathbf{s}$ 合并成单个矩阵：

\begin{equation}
\begin{pmatrix}
R & \mathbf{s}\\
\bzero^{\!\top} & 1
\end{pmatrix}
\end{equation}

这是一个 $5\times5$ 矩阵；我们用符号 $R$ 缩写整个 $4\times4$ 矩阵，用符号 $\mathbf{s}$ 表示矩阵最后一列的前四个分量，用符号 $\bzero$ 表示矩阵最后一行中的前四个分量。用末尾带坐标 1 的矢量去乘它，即得：

\begin{equation}
\begin{pmatrix}
R & \mathbf{s}\\
\bzero^{\!\top} & 1
\end{pmatrix}
\begin{pmatrix}
\bv\\
1
\end{pmatrix}
=
\begin{pmatrix}
R\bv + \mathbf{s}\\
1
\end{pmatrix}
\end{equation}

我们可以通过把两个这样的对称性的矩阵相乘，来看出它们相继施加时如何相互作用，其中第一个对称性的矩阵在矩阵乘积中位于右侧：

\begin{equation}
\begin{pmatrix}
R_{2} & \mathbf{s}_{2}\\
\bzero^{\!\top} & 1
\end{pmatrix}
\begin{pmatrix}
R_{1} & \mathbf{s}_{1}\\
\bzero^{\!\top} & 1
\end{pmatrix}
=
\begin{pmatrix}
R_{2}R_{1} & R_{2}\mathbf{s}_{1} + \mathbf{s}_{2}\\
\bzero^{\!\top} & 1
\end{pmatrix}
\end{equation}

\section{进阶：自由经典粒子的哈密顿量与拉格朗日量}

在关于能量-动量矢量的入门一节中，我们考虑了一个速度为 $v$ 的自由经典粒子，并推导出它的总能量 $E$ 与动量 $p$（定义为能量-动量矢量的 $t$ 分量与 $x$ 分量）的公式：

\begin{equation}
E = \frac{m}{\sqrt{1+v^{2}}} \tag{1}
\end{equation}

\begin{equation}
p = \frac{mv}{\sqrt{1+v^{2}}} \tag{2}
\end{equation}

\begin{equation}
E^{2} = m^{2} - p^{2} \tag{3}
\end{equation}

我们现在的目标是用哈密顿力学与拉格朗日力学来描述同一个粒子。这不会告诉我们关于该粒子行为的任何新东西——我们已经知道它只是沿欧几里得四维空间中的一条直线世界线运动——但它将使我们能够选对符号与约定，让一切都自洽地运作，这在我们想要把哈密顿与拉格朗日形式体系用于更精密的目的时会大有帮助。

首先，我们假定已选定一套正交坐标 $x$、$y$、$z$、$t$。我们再假定粒子只沿 $x$ 方向运动来简化分析，于是只需考虑一个空间坐标；推广到三个空间坐标是直截了当的。

在哈密顿力学中，哈密顿量 $H$ 是把系统的总能量表示为 $n$ 个广义坐标 $q^{i}$（$i=1,\dots,n$，代表系统的自由度）、与这些坐标\textbf{共轭}（conjugate）的 $n$ 个动量 $p_{i}$、以及时间 $t$ 的函数。哈密顿方程给出动量与坐标的变化率\footnote{见参考文献 [1]。}：

\begin{equation}
\frac{dp_{i}}{dt} = -\frac{\partial H}{\partial q^{i}} \tag{4}
\end{equation}

\begin{equation}
\frac{dq^{i}}{dt} = \frac{\partial H}{\partial p_{i}} \tag{5}
\end{equation}

对于我们这个沿 $x$ 方向运动的自由粒子，$x$ 是唯一的坐标，将有一个共轭于它的动量 $p$。于是乍看之下，我们或许会以为可以把方程 \textbf{(3)} 中的 $E$ 直接代入得到 $H(x,p,t)=E$：

\begin{equation}
H(x,p,t) = \sqrt{m^{2}-p^{2}} \tag{6}
\end{equation}

而共轭于坐标 $x$ 的动量就简单地是能量-动量矢量的 $x$ 分量。

然而，这种做法行不通！用这套设定，哈密顿第二方程 \textbf{(5)} 会给出：

\begin{equation}
\begin{aligned}
\frac{dx}{dt} &= \frac{\partial H}{\partial p}
= \frac{\partial\bigl[\sqrt{m^{2}-p^{2}}\bigr]}{\partial p}
= -\frac{p}{\sqrt{m^{2}-p^{2}}}
= -v \qquad \textbf{[错！]}
\end{aligned}
\tag{7}
\end{equation}

我们的错误在于，想当然地认为按几何方式定义的动量——即能量-动量矢量的 $x$ 分量——必定与哈密顿形式体系中共轭于 $x$ 的动量相同。

但这很容易修正。如果我们定义一个“哈密顿动量”：

\begin{equation}
p_{H} = -p = -\frac{mv}{\sqrt{1+v^{2}}} \tag{8}
\end{equation}

作为共轭于坐标 $x$ 的动量，那么一切就都漂亮地自洽了：

\begin{equation}
H(x,p_{H},t) = \sqrt{m^{2}-p_{H}^{2}} \tag{9}
\end{equation}

\begin{equation}
\begin{aligned}
\frac{dx}{dt} &= \frac{\partial H}{\partial p_{H}}
= \frac{\partial\bigl[\sqrt{m^{2}-p_{H}^{2}}\bigr]}{\partial p_{H}}
= -\frac{p_{H}}{\sqrt{m^{2}-p_{H}^{2}}}
= v \qquad \textbf{[对！]}
\end{aligned}
\tag{10}
\end{equation}

这看似微不足道，但在此处把符号搞对，比等到我们遇到黎曼热力学中温度可正可负、正确选择未必那么显然的种种奇特之处时再来收拾要容易得多。

哈密顿第一方程 \textbf{(4)} 给出：

\begin{equation}
\frac{dp_{H}}{dt} = -\frac{\partial H}{\partial x} = 0 \tag{11}
\end{equation}

它只不过告诉我们粒子的动量是常量。

接下来我们来看自由粒子的拉格朗日量。拉格朗日量 $L$ 是这样一个量：它关于时间的积分——称为作用量（action）$S$——在粒子轨迹的变分下取极大或极小\footnote{见参考文献 [2]。}。这立刻给了我们一条线索：如果我们令 $L\,dt$ 正比于粒子世界线在时间间隔 $dt$ 内那一段的长度，那么作用量 $S=\int L\,dt$ 将正比于世界线的长度。由于两点之间最短的路径是直线，极小化作用量将产生直线的世界线，这正是自由粒子应有的样子。

拉格朗日量 $L$ 形式上定义为系统广义坐标 $q^{i}$、这些坐标的时间变化率 $\dot{q}^{i}=dq^{i}/dt$、以及时间 $t$ 的函数。对本系统，只有一个坐标 $x$，其时间变化率为 $v$，我们有：

\begin{equation}
\begin{aligned}
L(x,v,t)\,dt &= C\,ds,\\
L(x,v,t) &= C\,\frac{ds}{dt},
\end{aligned}
\tag{12}
\end{equation}

其中 $s$ 是沿粒子世界线的四维空间距离，$C$ 是尚待确定的常数。

对于在我们选定坐标系中以速度 $v$ 运动的粒子，沿其世界线的四维空间距离 $s$ 直接由勾股定理给出：它的平方等于流逝的时间 $t$ 与通过空间所走过的距离 $x=vt$ 的平方之和：

\begin{equation}
\begin{aligned}
s(t) &= \sqrt{t^{2} + v^{2}t^{2}}\\
&= t\sqrt{1+v^{2}}
\end{aligned}
\tag{13}
\end{equation}

这与 \textbf{(12)} 一起给出：

\begin{equation}
\begin{aligned}
\frac{ds}{dt} &= \sqrt{1+v^{2}},\\
L(x,v,t) &= C\sqrt{1+v^{2}}
\end{aligned}
\tag{14}
\end{equation}

一般的欧拉-拉格朗日方程为\footnote{见参考文献 [2]。}：

\begin{equation}
\frac{d}{dt}\left[\frac{\partial L}{\partial \dot{q}^{i}}\right] = \frac{\partial L}{\partial q^{i}} \tag{15}
\end{equation}

对本系统它变为：

\begin{equation}
\begin{aligned}
\frac{d}{dt}\left[\frac{\partial L}{\partial v}\right] &= \frac{\partial L}{\partial x},\\
\frac{d}{dt}\left[\frac{Cv}{\sqrt{1+v^{2}}}\right] &= 0,\\
\left[\frac{C}{m}\right]\frac{dp}{dt} &= 0
\end{aligned}
\tag{16}
\end{equation}

于是无论我们为常数 $C$ 取什么值，这都只会再次告诉我们粒子的动量是常量。不过，我们可以利用哈密顿力学与拉格朗日力学之间的关系来确定 $C$ 的正确取值——在后者中，哈密顿共轭动量是用拉格朗日量定义的\footnote{见参考文献 [3]。}：

\begin{equation}
p_{i} = \frac{\partial L}{\partial \dot{q}^{i}} \tag{17}
\end{equation}

对本系统，利用 \textbf{(8)} 与 \textbf{(14)}，它变为：

\begin{equation}
\begin{aligned}
p_{H} &= \frac{\partial L}{\partial v},\\
-\frac{mv}{\sqrt{1+v^{2}}} &= \frac{Cv}{\sqrt{1+v^{2}}},\\
C &= -m.
\end{aligned}
\end{equation}

于是我们得到拉格朗日量的最终表达式：

\begin{equation}
L(x,v,t) = -m\sqrt{1+v^{2}} \tag{18}
\end{equation}

作为交叉验证，我们可以确认哈密顿量本身与拉格朗日量之间的关系正确无误，一般而言\footnote{见参考文献 [1]。}：

\begin{equation}
H = p_{i}\dot{q}^{i} - L \tag{19}
\end{equation}

对本系统它变为：

\begin{equation}
\begin{aligned}
H &= p_{H}v - L\\
&= v\left(-\frac{mv}{\sqrt{1+v^{2}}}\right) - \left(-m\sqrt{1+v^{2}}\right)\\
&= m\left[\sqrt{1+v^{2}} - \frac{v^{2}}{\sqrt{1+v^{2}}}\right]\\
&= \frac{m}{\sqrt{1+v^{2}}}
\end{aligned}
\tag{20}
\end{equation}

与方程 \textbf{(1)} 一致。

因此，要用哈密顿与拉格朗日形式体系处理黎曼物理，须注意：

\begin{itemize}
  \item 自由粒子的哈密顿量 $H$ 是其总能量，但要用一个“哈密顿动量” $p_{H}$ 表示，而它是通常的、几何定义的动量的\emph{相反数}；
  \item 质量为 $m$ 的自由粒子的作用量 $S$ 等于 $-m$ 乘以它世界线的普通欧几里得长度，而拉格朗日量 $L$ 是这一作用量关于时间的导数。
\end{itemize}

\begin{notebox}[补充：与洛伦兹情形的符号差异]
为便于对照（以下属于标准物理知识的补充说明）：在洛伦兹情形中，自由粒子相应的量是哈密顿量 $H=\sqrt{m^{2}+p^{2}}$（即 $E=\sqrt{m^{2}+p^{2}}$）、拉格朗日量 $L=-m\sqrt{1-v^{2}}$，且共轭动量就是几何动量本身。而在黎曼情形中，根号内 $p^{2}$ 前是\emph{减号}（$H=\sqrt{m^{2}-p^{2}}$），共轭动量是几何动量的\emph{反号}（$p_{H}=-p$），拉格朗日量则是 $L=-m\sqrt{1+v^{2}}$——三个符号差异恰好彼此协调，使哈密顿方程与欧拉-拉格朗日方程给出正确的 $dx/dt=v$。
\end{notebox}

\section{参考文献}

\begin{itemize}
  \item [1] L. D. Landau and E. M. Lifshitz, \emph{Mechanics}, Butterworth-Heinemann, 1976, 第 40 节。
  \item [2] Landau and Lifshitz, \emph{op. cit.}, 第 2 节。
  \item [3] Landau and Lifshitz, \emph{op. cit.}, 第 7 节。
\end{itemize}
